\begin{abstract}
Functional connectomes have an obvious graph structure, yet neighborhood
aggregation can harm prediction when every node already contains a global
connectivity profile. We examined whether learned orthogonal bundle transport
changes this density-dependent behavior. The analysis included
\StudyParticipants{} technically eligible ABIDE-I participants from
\StudySites{} sites. A single neural backbone was paired with identity
propagation, a learned-local capacity control, GCN aggregation, trivial-bundle
diffusion, or learned BuNN transport at positive edge densities of 0, 1, 5, 10,
and 20 percent. We used nested held-out-site validation, five final seeds,
regularized non-graph baselines, and representation diagnostics computed in a
common frame. The primary BuNN-minus-GCN performance-curve contrast was
\BunnGcnEstimate{} (95\% site-bootstrap CI [\BunnGcnCILow{},
\BunnGcnCIHigh{}]; exact sign-flip $p=\BunnGcnP$). Relative to the connectome
elastic net, the BuNN contrast was \BunnElasticEstimate{} (95\% CI
[\BunnElasticCILow{}, \BunnElasticCIHigh{}]). The matched-anchor
effective-rank contrast was \RankEstimate{} (95\% CI [\RankCILow{},
\RankCIHigh{}]), which ran counter to the proposed preservation advantage.
Site weighting and training seed affected the small BuNN-GCN difference. For
this ABIDE-I pipeline, learned bundle transport provided no detected
predictive or representation-preservation advantage. This finding concerns a
computational operator under a specified analysis; it does not establish a
biological bundle geometry or a general result for BuNNs.
\end{abstract}
